\subsection{Brennweite des Achromats}
Das eingangs durchgeführte Experiment mit Untersuchung der Abhängigkeiten von \(g\) und \(b\) zeigte, wie intuitiv aufgrund der Existenz von Lupen erwartet, dass mit kleinerer Gegenstandsweite der Abbildungsmaßstab vergrößert wird. Die Effekte, die auftreten, wenn \(g\leq f\) wurden mir experimentell während des Versuches nicht klar. Erst durch Studieren der physikalischen Grundlagen kam das Verständnis der auftretenden Effekte. Es ist zwar im Messprotokoll nicht vermerkt, aber dass man das Bild bei \(g=f\) mit größerem \(b\) immer besser zu sehen schien (\(b\rightarrow \infty\)), sowie die bei \(g<f\) nicht gefundene Scharfstellung, ist uns während des Versuches aufgefallen.

\subsection{Besselverfahren}
Die Herleitung der Besselformel \cref{eq:Bessel} wäre interessant zu sehen, dieses Verfahren hat man leider nur stur nach Skript durchgeführt. Da war die Bestimmung der Brennweite der achromatischen Linse in Aufgabe 1 intuitiver. Bei der Auswertung der Ergebnisse für die Brennweiten bei unterschiedlichen Wellenlängen wurde wie in der Theorie gelernt, die Dispersion deutlich. Im Skript\autocite{Skript_1} ist auch eine Formel zu finden, die die Abhängigkeit der Brennweite vom Brechungsindex \(n\) und den Krümmungsradien \(r_1,\;r_2\) mathematisch ausdrückt: 
\begin{equation}
    \frac{1}{f}=(n-1)\left(\frac{1}{r_1}+\frac{1}{r_2}\right)
    \label{N}
\end{equation}
Betrachtet man nun die Kurve des Brechungsindexes \(n(\lambda)\), der für transparente Materialen bei niedrigeren Wellenlängen größer wird (\autocite{Brechungsindex}), so folgt daraus, dass nach \(f\sim \frac{1}{n}\) die Brennweite für kleinere Wellenlängen wie von rotem oder blauen Licht sinkt. Diese Überlegung stimmt auch mit den experimentell nach der Besselformel bestimmten Werten überein. Die Messung von \(f_{\text{rot}}\) weicht um \(\qty{+0,2}{\cm}=\qty{+1,64}{\%}\) von \(f_{\text{blau}}\) ab. Verglichen mit der prozentualen Messunsicherheit von \qty{0,58}{\%} ist das eine aussagekräftige Messung im Sinne des eben beschriebenen Dispersionseffektes.  

\subsection{Linsenfeherquellen}
Die bei der bikonvexen Linse entstehende chromatische Abberation haben wir nicht genau erkennen können, jedoch gab es auch keine Ambition, sie näher zu untersuchen. Die sphärische Abberation hingegen war durch das Verwenden der Lochblende sehr gut zu erkennen. Interessant wäre, den Versuch dahingehend zu erweitern, ob durch versetzt stehende Gegenstände eine qualitative Untersuchung der Schärfentiefe möglich ist. Die fehlende Schärfe der Ringblende verglichen mit der Lochblende geht einher mit den eingangs besprochenen physikalischen Ursachen der sphärischen Abberation. Die Ringblende beschränkt sich auf achsenfernere Strahlen, die die sphärische Abberation aufgrund der nicht optimalen Geometrie der Linse deutlich erhöhen.

\subsection{Mikroskopauflösung}
Der sehr hohe Fehler des über den Öffnungswinkel bestimmten Wertes für \(G_{\text{min}}\) hat seine Ursache in dem exorbitanten relativen Fehler \(\Delta d=\qty{50}{\%}\). Das liegt daran, dass die Skala nur unter Einfluss von Parallaxeneffekten abgelesen werden konnte und die mechanische Einstellung der Spaltbreite nicht sehr genau war. 

\subsection{Fazit}
Zwar wurden einige interessante Experimente durchgeführt, um ein besseres Gefühl für optische Systeme zu bekommen, ein super intuitives Verständnis, das ermöglichen würde, selber Linsenanordnungen zu entwickeln, die bestimmte Zwecke erfüllen sollen, habe ich leider nicht erlangt. Der exakte Einfluss des Auges in optischen Systemen wurde nur durch das sehr zeitaufwändige Studieren von anderer Literatur halbwegs verständlich, zumindest in den für hier erforderlichen Ausmaßen. Zusätzlich fehlt auch ein wenig die tiefergehende Betrachtung des Auges und seiner Funktionsweise. Dennoch wurden ein paar schöne Dinge bestimmt und in der Auswertung deutlich klarer.
